\documentclass[10pt,a4paper]{amsart}
\usepackage{amsmath,amssymb,amsthm,graphicx}
% Retained minimal preamble for the Stacks Project chapter more-algebra.tex (More on Algebra).
% The chapter file pinned-original-source.tex carries only \input{preamble} on its first line, so no
% preamble text lives inside the chapter itself.  The Stacks Project uses one single shared preamble for
% every chapter; the snapshot stages it identically as sources/stacks-chow/preamble.tex and
% sources/stacks-etale-cohomology/preamble.tex.  The block below is copied mechanically, line for line,
% from sources/stacks-chow/preamble.tex lines 14-18 (the Xy-pic setup) and 154-217 (the native theorem
% declarations together with the mathematical macros).  Nothing was added, removed or redefined; the
% wrapper class and the amsmath/amssymb/amsthm/graphicx base packages are supplied by the assembler.
% Chapter: more-algebra.tex -- More on Algebra.
% For commutative diagrams we use Xy-pic
\usepackage[all]{xy}

% We use 2cell for 2-commutative diagrams.
\xyoption{2cell}

% Theorem environments.
%
\theoremstyle{plain}
\newtheorem{theorem}[subsection]{Theorem}
\newtheorem{proposition}[subsection]{Proposition}
\newtheorem{lemma}[subsection]{Lemma}

\theoremstyle{definition}
\newtheorem{definition}[subsection]{Definition}
\newtheorem{example}[subsection]{Example}
\newtheorem{exercise}[subsection]{Exercise}
\newtheorem{situation}[subsection]{Situation}

\theoremstyle{remark}
\newtheorem{remark}[subsection]{Remark}
\newtheorem{remarks}[subsection]{Remarks}

\numberwithin{equation}{subsection}

% Macros
%
\def\lim{\mathop{\mathrm{lim}}\nolimits}
\def\colim{\mathop{\mathrm{colim}}\nolimits}
\def\Spec{\mathop{\mathrm{Spec}}}
\def\Hom{\mathop{\mathrm{Hom}}\nolimits}
\def\Ext{\mathop{\mathrm{Ext}}\nolimits}
\def\SheafHom{\mathop{\mathcal{H}\!\mathit{om}}\nolimits}
\def\SheafExt{\mathop{\mathcal{E}\!\mathit{xt}}\nolimits}
\def\Sch{\mathit{Sch}}
\def\Mor{\mathop{\mathrm{Mor}}\nolimits}
\def\Ob{\mathop{\mathrm{Ob}}\nolimits}
\def\Sh{\mathop{\mathit{Sh}}\nolimits}
\def\NL{\mathop{N\!L}\nolimits}
\def\CH{\mathop{\mathrm{CH}}\nolimits}
\def\proetale{{pro\text{-}\acute{e}tale}}
\def\etale{{\acute{e}tale}}
\def\QCoh{\mathit{QCoh}}
\def\Ker{\mathop{\mathrm{Ker}}}
\def\Im{\mathop{\mathrm{Im}}}
\def\Coker{\mathop{\mathrm{Coker}}}
\def\Coim{\mathop{\mathrm{Coim}}}

% Boxtimes
%
\DeclareMathSymbol{\boxtimes}{\mathbin}{AMSa}{"02}

%
% Macros for moduli stacks/spaces
%
\def\QCohstack{\mathcal{QC}\!\mathit{oh}}
\def\Cohstack{\mathcal{C}\!\mathit{oh}}
\def\Spacesstack{\mathcal{S}\!\mathit{paces}}
\def\Quotfunctor{\mathrm{Quot}}
\def\Hilbfunctor{\mathrm{Hilb}}
\def\Curvesstack{\mathcal{C}\!\mathit{urves}}
\def\Polarizedstack{\mathcal{P}\!\mathit{olarized}}
\def\Complexesstack{\mathcal{C}\!\mathit{omplexes}}
% \Pic is the operator that assigns to X its picard group, usage \Pic(X)
% \Picardstack_{X/B} denotes the Picard stack of X over B
% \Picardfunctor_{X/B} denotes the Picard functor of X over B
\def\Pic{\mathop{\mathrm{Pic}}\nolimits}
\def\Picardstack{\mathcal{P}\!\mathit{ic}}
\def\Picardfunctor{\mathrm{Pic}}
\def\Deformationcategory{\mathcal{D}\!\mathit{ef}}

% --- appended 2026-10-12 (Chief): the project own definitions of the reference/slogan/history
% display environments, copied mechanically from the shared Stacks preamble lines 8-12
% (\usepackage{verbatim} plus the three \newenvironment lines, which the chapter preamble keeps
% verbatim from the project preamble). Nothing added, removed or redefined: these environments
% carry no mathematical content and the project itself typesets them as comments. Without these
% three lines a counted statement whose own text uses them fails with Environment slogan undefined.
\usepackage{verbatim}
\newenvironment{reference}{\comment}{\endcomment}
\newenvironment{slogan}{\comment}{\endcomment}
\newenvironment{history}{\comment}{\endcomment}
\usepackage{cleveref}
\crefname{originallemma}{original lemma}{original lemmas}
\Crefname{originallemma}{Original lemma}{Original lemmas}
\makeatletter
\newlabel{lemma-choose-values}{{0GVA}{source}{}{sourcepointer-1}{}}
\newlabel{lemma-choose-values@cref}{{[originallemma][0][]0GVA}{[1][0][]source}{}{}{}}
\makeatother
\crefname{originallemma}{original lemma}{original lemmas}
\Crefname{originallemma}{Original lemma}{Original lemmas}
\makeatletter
\newlabel{algebra-lemma-chinese-remainder}{{00DT}{source}{}{sourcepointer-2}{}}
\newlabel{algebra-lemma-chinese-remainder@cref}{{[originallemma][0][]00DT}{[1][0][]source}{}{}{}}
\makeatother
\crefname{originalsection}{original section}{original sections}
\Crefname{originalsection}{Original section}{Original sections}
\makeatletter
\newlabel{topology-section-noetherian}{{0050}{source}{}{sourcepointer-3}{}}
\newlabel{topology-section-noetherian@cref}{{[originalsection][0][]0050}{[1][0][]source}{}{}{}}
\makeatother
\crefname{originalsection}{original section}{original sections}
\Crefname{originalsection}{Original section}{Original sections}
\makeatletter
\newlabel{topology-section-catenary-spaces}{{02I0}{source}{}{sourcepointer-4}{}}
\newlabel{topology-section-catenary-spaces@cref}{{[originalsection][0][]02I0}{[1][0][]source}{}{}{}}
\makeatother
\title{Splitting off the vanishing locus of one section while keeping the codimension bound}
\author{The Stacks Project Authors}
\date{stacks-project:more-algebra.tex:local-snapshot-20261012; GNU FDL 1.2 or later}
\begin{document}
\maketitle
\setcounter{subsection}{1}
\noindent\textit{Editorial note.} All delivered mathematical text is the author's own, copied from the pinned Stacks Project chapter file more-algebra.tex; no statement or proof step was rewritten, reordered or removed. The counted span is the author environment and its own proof environment. References inside the span are handed to the assembler as navigation stubs that carry the canonical tag of the cited tag only; whether each of them is load-bearing is recorded on the stub itself.\par\medskip
\noindent\textit{Allowed standard prerequisites (editorial).}\begin{itemize}
\item The situation fixing the sheaf of modules, the finite dimensional spaces of values and the sections, whose fields and dimension function the statement uses. used from the statement at line 38811 onward. Fixes the notation M, V(x), s\_1 to s\_h and F so that the statement and the proof talk about the same objects.
\item For the case of the base value of the induction one can prescribe the values of a section at finitely many points. invoked at proof lines 38835-38837. Produces the section s and the choice F' equal to the whole space in the case k = 0.
\item A finite set of functions can be prescribed at one finite set of points and made to vanish at another, by the Chinese remainder argument for the ring of functions. invoked at proof lines 38866-38868. Produces the element f equal to one at the y\_j and zero at the z\_l, which builds the perturbed section u plus f t.
\item The codimension function on a Noetherian topological space and the finiteness of the number of irreducible components. used throughout the proof from lines 38831-38835 and at many later steps. Lets the proof decompose the closed sets G and H into irreducible components, pick points outside the other components, and compare codimensions.
\end{itemize}\medskip
\ifdefined\hypertarget\hypertarget{sourcepointer-1}{}\fi
\noindent\textit{Source reference.} Original lemma 0GVA: sources/stacks-more-algebra/pinned-original-source.tex, source lines 38795--38795. Used at proof lines 38835-38837 for the base case of the induction, where the whole space is taken as the closed set so no codimension estimate is needed; every other step of the induction is inside the counted span.\par\medskip
\ifdefined\hypertarget\hypertarget{sourcepointer-2}{}\fi
\noindent\textit{Source reference.} Original lemma 00DT: sources/stacks-chow/tags.txt, source lines 1--1. Invoked at proof lines 38866-38868 to produce the element f taking value one at the chosen points of the first closed set and zero at those of the second; the construction of the perturbed section and its verification are inside the counted span.\par\medskip
\ifdefined\hypertarget\hypertarget{sourcepointer-3}{}\fi
\noindent\textit{Source reference.} Original section 0050: sources/stacks-chow/tags.txt, source lines 1--1. Cited at proof lines 38832-38835 as the source of the finiteness and decomposition statements for closed sets; the decompositions into irreducible components used later are performed inside the counted span.\par\medskip
\ifdefined\hypertarget\hypertarget{sourcepointer-4}{}\fi
\noindent\textit{Source reference.} Original section 02I0: sources/stacks-chow/tags.txt, source lines 1--1. Cited at proof lines 38831-38833 to point at the definition of codimension used in the statement and in the comparisons; all codimension comparisons are made inside the counted span.\par\medskip
\begin{proposition}
\label{proposition-splitting}
\begin{reference}
\cite[Theorem 2]{Serre-projective}
\end{reference}
In Situation \ref{situation-splitting} assume $\Omega$ is
a Noetherian topological space. Let $s_1, \ldots, s_h \in M$.
Let $Z(s_1, \ldots, s_h) \subset F \subset \Omega$ be closed.
Let $x_1, \ldots, x_n \in F$ be pairwise distinct.
Let $v_i  \in V(x_i)$.
Let $k \geq 0$ be an integer such that
$$
(*)\quad h + k \leq \dim_{\kappa(x)} V(x)\text{ for all }x \in \Omega
$$
Then there exist $s \in M$ and $F' \subset \Omega$ closed such that
\begin{enumerate}
\item[(a)] $s(x_i)$ maps to $v_i$,
\item[(b)] $Z(s_1, \ldots, s_h, s) \subset F \cup F'$, and
\item[(c)] every irreducible component of $F'$ has
codimension $\geq k$ in $\Omega$.
\end{enumerate}
\end{proposition}

\begin{situation}
\label{situation-splitting}
Here $R$ is a ring and $M$ is a finitely presented $R$-module.
Denote $\Omega \subset \Spec(R)$ the set of closed points
with the induced topology. For $x \in \Omega$ denote
$M(x) = M/xM$ the fibre of $M$ at $x$. This is a finite
dimensional vector space over the residue field $\kappa(x)$ at $x$.
Given $s \in M$ we denote $s(x)$ the image of $s$ in $M(x)$.
\end{situation}

\begin{proof}
We note that codimension was defined in
Topology, Section \ref{topology-section-catenary-spaces}
and that we will use some results on Noetherian topological
spaces contained in
Topology, Section \ref{topology-section-noetherian}.

\medskip\noindent
The proof is by induction on $k$. If $k = 0$, then we choose
$s \in M$ as in Lemma \ref{lemma-choose-values} and we choose
$F' = \Omega$.

\medskip\noindent
Assume $k > 0$. By our induction hypothesis we may
choose $u \in M$ and $G \subset \Omega$ closed satisfying
(a), (b), (c) for $s_1, \ldots, s_h$, $F$, $x_1, \ldots, x_n$,
$v_1, \ldots, v_n$, and $k - 1$.

\medskip\noindent
Let $G = G_1 \cup \ldots \cup G_m$ be the decomposition of $G$ into
its irreducible components. If $G_j \subset F$, then we can remove
it from the list. Thus we may assume $G_j$ is not contained in $F$
for $j = 1, \ldots, m$.
For $j = 1, \ldots, m$ choose $y_j \in G_j$ with $y_j \not \in F$
and $y_j \not \in G_{j'}$ for $j' \not = j$. This is possible as
there are no inclusions among the irreducible components of $G$.
Choose $w_j \in V(y_j)$ not contained in the span of the images of
$s_1(y_j), \ldots, s_h(y_j)$; this is possible because
$h + k \leq \dim V(y_j)$ and $k > 0$.

\medskip\noindent
Apply the induction hypothesis to the $h + 1$ sections
$s_1, \ldots, s_h, u$, the closed set $F \cup G$,
the points $x_1, \ldots, x_n, y_1, \ldots, y_m \in F \cup G$,
the elements $0 \in V(x_i)$ and $w_j \in V(y_j)$, and
the integer $k - 1$. Note that we have increased $h$ by $1$ and
decreased $k$ by $1$ hence the assumption $(*)$ of the proposition
remains valid. This produces $t \in M$ and $H \subset \Omega$ closed
satisfying (a), (b), (c) for $s_1, \ldots, s_h, u$, $F \cup G$,
$x_1, \ldots, x_n, y_1, \ldots, y_m$, $0, \ldots, 0, w_1, \ldots, w_m$,
and $k - 1$.

\medskip\noindent
Let $H_1, \ldots, H_p \subset H$ be the irreducible components of $H$
which are not contained in $F \cup G$. As before pick $z_l \in H_l$,
$z_l \not \in F \cup G$ and $z_l \not \in H_{l'}$ for $l' \not = l$.
Using Algebra, Lemma \ref{algebra-lemma-chinese-remainder}
we may choose $f \in R$ such that
$f(y_j) = 1$, $j = 1, \ldots, m$ and $f(z_l) = 0$, $l = 1, \ldots, p$.
Claim: the element $s = u + f t$ works.

\medskip\noindent
First, the value $s(x_i)$ agrees with $u(x_i)$ because
$t(x_i) = 0$ and hence we see that $s(x_i)$ maps to $v_i$.
This proves (a).
To finish the proof it suffices to show that every irreducible
component $Z$ of $Z(s_1, \ldots, s_h, s)$ not contained in $F$
has codimension $\geq k$ in $\Omega$. Namely, then we can set $F'$
equal to the union of these and we get (b) and (c).
We can see that irreducible components $Z$ of $Z(s_1, \ldots, s_h, s)$
of codimension $\leq k - 1$ do not exist as follows:
\begin{enumerate}
\item Observe that
$Z(s_1, \ldots, s_h, s) \subset Z(s_1, \ldots, s_h, u, t) = F \cup H$
as $s = u + ft$. Hence $Z \subset H$.
\item The irreducible components of $H$ have codimension $\geq k - 1$.
Hence $Z$ is equal to an irreducible component of $H$
as $Z$ has codimension $\leq k - 1$.
Hence $Z = H_l$ for some $l \in \{1, \ldots, p\}$ or
$Z = G_j$ for some $j \in \{1, \ldots m\}$.
\item But $Z = G_j$ is impossible as $s_1(y_j), \ldots, s_h(y_j)$
map to linearly independent elements of $V(y_j)$ and
$s(y_j) = u(y_j) + f(y_j) t(y_j) = u(y_j) + t(y_j)$
maps to an element of the form
$$
\text{linear combination images of }s_i(y_j) + w_j
$$
which is linearly independent of the images of $s_1(y_j), \ldots, s_h(y_j)$
in $V(y_j)$ by our choice of $w_j$.
\item Also $Z = H_l$ is impossible. Namely, again
$s_1(z_l), \ldots, s_h(z_l)$ map to linearly independent
elements of $V(z_l)$ and $s(z_l) = u(z_l) + f(z_l) t(z_l) = u(z_l)$
maps to an element of $V(z_l)$ linearly independent of those
as $z_l \not \in F \cup G$.
\end{enumerate}
This finishes the proof.
\end{proof}

\end{document}
